Training data contamination inflates Large Language Model (LLM) benchmark scores, yet no existing method provides uncertainty-quantified calibration grounded in psychometric measurement theory. We introduce ContamCal, a two-stage framework that fuses multiple detection signals via training-free Reciprocal Rank Fusion (RRF) into continuous per-item exposure scores, then embeds these scores in a contamination-augmented Item Response Theory (IRT) model to decompose observed performance into latent ability and contamination effects. Applied to MMLU, ARC-Challenge, and GSM8K with controlled contamination injection across the Qwen2.5 model family ($J=13$), we find that IRT ability estimation alone reduces calibration error by up to 53\% relative to raw scores, independent of whether contamination is explicitly modeled. The contamination parameter $\gamma$ provides complementary diagnostic information: exploratory per-domain analysis reveals directional heterogeneity across MMLU cognitive domains (positive $\gamma$ in STEM, negative in Social Sciences), though most estimates are underpowered at this ensemble size. We further document that Min-K\%++ exhibits directional reversal at small model scales (Cohen's $d = -1.55$ to $-1.98$), producing actively misleading conclusions. Simulation studies establish that reliable $\gamma$ recovery requires signal AUC $\geq 0.55$ and ensemble size $J \geq 50$, quantifying a structural gap between current LLM evaluation practice and the statistical requirements of contamination-augmented IRT.